\chapter{The Editors}\label{cha:editors}

Two of them, because they answer different questions. \verb|EDIT| is the
one to reach for when a file wants a line changed and you want to be out
again in ten seconds. \verb|VI| is the one for a file too big to think
about, and it is modal, so it expects you to have met \texttt{vi} before.

Both draw through JIM, the VT100 in hardware (Chapter~\ref{cha:tube}), and
both take their keys raw from the ROM --- an arrow arrives as one byte
rather than an escape sequence to unpick, with a bit beside it saying
that it \emph{is} an arrow and not the accented letter that shares its
code. Neither needs the Tube.

\section{EDIT}
\verb|EDIT name| opens a file, or starts an empty one if there is no such
file yet. The whole text sits in memory below the program, so it holds
about 22\,KB --- ample for a \texttt{STARTUP.BAT}, a \texttt{.SUB}, a
BASIC listing or a letter. There are no modes: what you type goes in,
and everything else is a control key.

\begin{center}
\small
\begin{tabular}{@{}l >{\raggedright\arraybackslash}p{0.62\linewidth}@{}}
\toprule
\textbf{Key} & \textbf{Does} \\
\midrule
arrows & move by a character or a line \\
Home, End & start and end of the line \\
PgUp, PgDn & a screenful \\
Enter & split the line at the cursor \\
Backspace & rub out to the left; at the start of a line, join it to the one above \\
Delete & rub out under the cursor \\
Ctrl-O \emph{or} Ctrl-S & save \\
Ctrl-X & leave --- twice on a changed file, to throw the changes away \\
Ctrl-R & renumber a BASIC program: 10, 20, 30\ldots, and every \texttt{GOTO} with it \\
\bottomrule
\end{tabular}
\end{center}

Accented letters type as themselves. The bar along the bottom carries
the name, a \verb|*| while there are
unsaved changes, where the cursor is, and the keys to save and leave ---
and Ctrl-R too, on a BASIC file --- so there is nothing to remember. Long
lines slide sideways rather than wrap.

\section{VI}
\verb|VI name| is modal in the old way. In \emph{normal} mode the keys
are commands; a handful of them put you into \emph{insert} mode, where
what you type goes into the file; and Escape brings you back. The
\verb|:| line at the bottom is the third place, for the ex commands. The
cursor says which of the three you are in: a block in normal mode, a bar
while inserting, an underline on the \verb|:| line. Accented letters go
in like any other, in insert mode and after \verb|r|.

\subsection*{Modes}
\begin{center}
\small
\begin{tabular}{@{}l >{\raggedright\arraybackslash}p{0.62\linewidth}@{}}
\toprule
\textbf{Key} & \textbf{Enters insert\ldots} \\
\midrule
\texttt{i} & before the cursor \\
\texttt{a} & after the cursor \\
\texttt{I} & at the start of the line \\
\texttt{A} & at the end of the line \\
\texttt{o} & on a new line below \\
\texttt{O} & on a new line above \\
\texttt{s} & replacing the character under the cursor \\
\texttt{S} & replacing the whole line \\
\texttt{C} & replacing the rest of the line \\
\texttt{c}\emph{motion} & replacing what the motion covers; \texttt{cc} the whole line \\
\midrule
Esc & \ldots and leaves it, back to normal mode \\
\texttt{:} & opens the command line (from normal mode) \\
\bottomrule
\end{tabular}
\end{center}

\subsection*{Moving}
A count goes in front of nearly anything: \verb|5j| is five lines down,
\verb|10G| is line ten.

\begin{center}
\small
\begin{tabular}{@{}l >{\raggedright\arraybackslash}p{0.62\linewidth}@{}}
\toprule
\textbf{Key} & \textbf{Moves} \\
\midrule
\texttt{h} \texttt{j} \texttt{k} \texttt{l} & left, down, up, right --- the arrows do the same \\
\texttt{w} & to the next word \\
\texttt{b} & back a word \\
\texttt{e} & to the end of the word \\
\texttt{0} & to the start of the line \\
\texttt{\textasciicircum} & to the first non-blank on the line \\
\texttt{\$} & to the end of the line \\
\texttt{gg} & to the first line \\
\texttt{G} & to the last line, or to line \emph{n} with a count \\
PgUp, PgDn & a screenful \\
\texttt{/}\emph{text} & to the next \emph{text} --- a plain substring, not a pattern \\
\texttt{?}\emph{text} & to the previous one \\
\texttt{n} & the same search again; \texttt{N} the other way \\
\bottomrule
\end{tabular}
\end{center}

\subsection*{Changing}
\begin{center}
\small
\begin{tabular}{@{}l >{\raggedright\arraybackslash}p{0.62\linewidth}@{}}
\toprule
\textbf{Key} & \textbf{Does} \\
\midrule
\texttt{x} & delete the character under the cursor \\
\texttt{X} & delete the one to its left \\
\texttt{d}\emph{motion} & delete what the motion covers: \texttt{dw} a word, \texttt{d\$} to the end \\
\texttt{dd} & delete the line \\
\texttt{D} & delete to the end of the line \\
\texttt{y}\emph{motion}, \texttt{yy} & copy (yank) the same, without deleting \\
\texttt{p} & put what was deleted or yanked, after the cursor \\
\texttt{P} & put it before \\
\texttt{r}\emph{c} & replace one character with \emph{c} \\
\texttt{\textasciitilde} & flip the case of one character \\
\texttt{J} & join the next line onto this one \\
\texttt{u} & undo --- unlimited \\
Ctrl-R & redo --- also unlimited \\
\bottomrule
\end{tabular}
\end{center}

Operators take any motion --- \verb|d2w| deletes two words, \verb|y$| yanks
to the end of the line --- or double themselves to work on whole lines.
A charwise operator \textbf{clamps to the line}: \verb|dw| at the end of
a line stops there instead of eating into the next one. That is
deliberate, and so is the plain-substring search: \verb|/foo| finds
\verb|foo|, and \verb|.| \verb|*| \verb|[| mean themselves.

\subsection*{The command line}
\begin{center}
\small
\begin{tabular}{@{}l >{\raggedright\arraybackslash}p{0.62\linewidth}@{}}
\toprule
\textbf{Command} & \textbf{Does} \\
\midrule
\texttt{:w} & save \\
\texttt{:w} \emph{name} & save under another name \\
\texttt{:q} & leave --- refuses on a changed file \\
\texttt{:q!} & leave anyway, changes lost \\
\texttt{:wq}, \texttt{:x} & save and leave \\
\texttt{:s/}\emph{old}\texttt{/}\emph{new}\texttt{/} & replace the first \emph{old} on this line (\texttt{/g}: every one) \\
\texttt{:\%s/}\emph{old}\texttt{/}\emph{new}\texttt{/g} & the same on every line of the file \\
\texttt{:map} \emph{keys} \emph{result} & define a key sequence, in normal mode \\
\texttt{:imap} \emph{keys} \emph{result} & the same, in insert mode \\
\texttt{:renum} [\emph{start} [\emph{step}]] & renumber a BASIC program (10 and 10 unless told), and every \texttt{GOTO} with it; \texttt{u} puts it back \\
\texttt{:make} & save, compile a \texttt{.C} or \texttt{.PAS}, and go to the first error \\
\texttt{:run} & \texttt{:make}, and if it compiled, run it; a key comes back \\
\texttt{:cn}, \texttt{:cp} & the next, the previous compiler message \\
\texttt{:cc} \emph{n}, \texttt{:cl} & message \emph{n}; the whole list \\
\texttt{:set wrap} & a line longer than the screen is folded onto the rows under it (so it starts); \texttt{:set nowrap} keeps it to one row that scrolls sideways, and \texttt{:set wrap!} changes it over. Folded, \texttt{j} and \texttt{k} still move by a line of the file, and \texttt{gj} and \texttt{gk} by a row of the screen. \texttt{set nowrap} in VI.RC keeps it \\
\texttt{:set ts=}\emph{n} & Tab's width: it puts spaces to the next stop, never a tab character (4 unless told; \texttt{set ts=2} in VI.RC keeps it, for PROG too) \\
\bottomrule
\end{tabular}
\end{center}

\noindent The command word is read in either case, so \verb|:Q| and
\verb|:WQ| work with the caps lock on --- which is how most people
arrive in VI from a BASIC.

\subsection*{Compiling from VI}
\verb|:make| is the edit--compile loop without leaving the editor. It saves
the file, compiles it with the machine's own \verb|CC| or \verb|PAS|
(Chapter~\ref{cha:linux}) --- which one, the name decides --- and puts
the cursor on the first error, with the message on the status line:

\begin{type}
VI HELLO.C
:make
\end{type}

\noindent \texttt{error 1 of 3: ';' expected} --- and \verb|:cn| is the
next, \verb|:cp| the one before, \verb|:cl| all of them on one screen,
\verb|:cc 3| the third. Mad Pascal gives the column as well as the line,
and the cursor goes to it. An error in another file --- a unit, or
\texttt{k4510.h} --- is shown with that file's name and does not move the
cursor. A clean compile says what it made, \texttt{hello.prg: 2140 bytes},
and counts the warnings, which \verb|:cn| walks through the same way.

\verb|:run| is \verb|:make| and then the program, as if typed at the
prompt. It runs over VI by way of \verb|SWAP|, so VI is there again when
it ends; the program's last screen stays up until you press a key. VI
reads the file back afterwards, because a program may use the far memory
VI keeps the text in --- which is why \verb|:run| saves first, and why the
undo history starts again. From a VI that was itself started by
\verb|SWAP| (from a BASIC's \verb|*VI|, or RANGER), \verb|:run| cannot
nest and says so: leave VI and type the program's name.

The compilers write what they said to \pth{/SYSTEM/LOG/MAKE.ERR}, one
message a line in one form for every language (\texttt{FILE:LINE:COL:KIND:TEXT}),
and that file is all VI reads --- so a compiler added later needs a
script, not a new VI.

\subsection*{Making \texttt{jk} leave insert mode}
The usual mapping is the one most people want first, and it is typed at
the \verb|:| line exactly like this, with the Escape spelt out in angle
brackets:

\begin{type}
:imap jk <Esc>
\end{type}

\noindent From then on, typing \verb|j| then \verb|k| in insert mode
leaves it, and the two letters never reach the file. A lone \verb|j|
still does: the editor holds it for one second waiting for the
\verb|k| --- vim's own \texttt{timeoutlen}, and long enough for a
deliberate ``j, k'' --- and hands it over when nothing comes. \verb|<CR>| is the other
spelt-out key; everything else in the result is literal.

To have it every time, put the same line, without the colon, in
\pth{/SYSTEM/ETC/VI.RC}: \verb|VI| runs that file once the text is
loaded, one ex command per line, a line starting with \verb|"| being a
comment. \pth{/SYSTEM/ETC/VI.SAMPLE} is one to copy, the way
\texttt{STARTUP.SAMPLE} is, and the file is yours --- it is not in the
repository. If a mapping you have defined seems not to fire, the usual
reason is that there is no \texttt{VI.RC} on that machine, only the
sample.

\begin{aside}
\textbf{Where the file lives.} Nothing of it is in the 64\,KB. Every line
is a 256-byte slot in far memory --- a length and up to 255 characters ---
and only the line under the cursor is brought down, fetched when the
cursor arrives and written back when it leaves. So the ceiling is far
memory rather than the address space: 32000 lines, and \verb|LOAD| and
\verb|SAVE| go straight to and from a flat copy out there without the file
ever passing through the CPU's view. Undo is a journal in the same
place, which is why it is unlimited: one level would have cost the same
as all of them.

A 1.28\,MB file of 20000 lines --- twenty times the whole address space
--- opens, edits at the far end and saves back byte for byte.
\end{aside}

\section{PROG}\label{sec:prog}
\verb|PROG name| is the front end for writing a program in C, Pascal or REXX:
the text, a menu bar, and --- under the text --- what the compiler said
about it. It is VI's engine with modern keys on it, in the spirit of
Turbo Pascal: F9 compiles, Ctrl-F9 compiles and runs, and an error puts
the cursor on the line it is about.

\begin{center}
\small
\begin{tabular}{@{}>{\raggedright\arraybackslash}p{0.33\linewidth} >{\raggedright\arraybackslash}p{0.62\linewidth}@{}}
\toprule
\textbf{Key} & \textbf{Does} \\
\midrule
arrows, Home, End, PgUp, PgDn & move; with Ctrl, a word at a time and the ends of the file \\
Enter, Tab & a new line that keeps the indent; spaces to the next tab stop (four, or \verb|set ts=| in VI.RC) --- with lines selected, Tab and Shift-Tab indent them and take the indent back \\
Shift with a moving key, Ctrl-A & select, by characters, from where the cursor was; Ctrl-A takes the whole file. Typing, Enter, Backspace and Delete replace what is selected \\
the mouse & a click places the cursor, opens a menu, brings a tab forward or goes to a message; a drag selects (Shift-click from the cursor); the wheel scrolls \\
Insert & insert or overwrite (the cursor is a bar or a block) \\
Ctrl-S \emph{or} F2, Ctrl-O, Ctrl-Q & save, open (in a tab of its own), quit --- asking first about every file with unsaved changes \\
Ctrl-N, Ctrl-W, F6, Shift-F6 & a new file, close this one, the next and the previous file \\
\emph{File} $\rightarrow$ \emph{New project} & a folder with a \texttt{PROJECT.K4P} and a first file that compiles as it stands \\
Ctrl-Z, Ctrl-Y & undo, redo, as far back as the session goes \\
Ctrl-X, Ctrl-C, Ctrl-V & cut, copy, paste --- the selection, or with nothing selected the line \\
Ctrl-F, F3, Ctrl-R, Ctrl-G & find, find again, replace everywhere, go to a line \\
Shift-Ctrl-F & find in files: every source file in this file's directory, into the message list \\
F9, Ctrl-F9 & save every changed file, compile the \texttt{.C} (\verb|CC|) or \texttt{.PAS} (\verb|PAS|); and run it. A \texttt{.RX} has nothing to compile: F9 saves it, Ctrl-F9 runs it, and an error takes you to its line \\
F4, Shift-F4 & the next, the previous message --- one about another file opens it \\
F10, F1 & the menu; the keys \\
\bottomrule
\end{tabular}
\end{center}

The messages are the ones \verb|:make| reads in VI, from the same
\pth{/SYSTEM/LOG/MAKE.ERR}: an error in the file on screen takes the
cursor to its line (and column, from Mad Pascal); one in another file ---
a unit, a header --- is listed with that file's name. A \verb|}| typed on
a line of its own goes back a level. F12 stays the machine's (its menu;
Shift+F12 pauses), so PROG leaves it alone; Ctrl-H is Backspace on this
keyboard, which is why replace is Ctrl-R.

Up to eight files are open at once, in the row under the menu bar --- the
one in front lit, a \verb|*| on each with unsaved changes. F9 saves every
changed file before it compiles, because the compilers read the disk: a
unit edited in another tab is what \verb|PAS| sees. An error in another
file --- a unit, a header --- opens that file (or brings its tab forward)
at the line. Find in files (Shift-Ctrl-F, or \emph{Search}) looks through
every source file in the directory of the file in front, and what it finds
goes into the message list, where F4 walks it like errors. After Ctrl-F9
the other open files are read back from the disk: a program may use the
memory they wait in.

\subsection*{Projects}
A program in several files is a \emph{project}: a folder, and in it a
\texttt{PROJECT.K4P} that names them, one \texttt{KEY=value} to a line:

\begin{type}
# GAME -- a PROG project
NAME=GAME
LANG=C
SRC=GAME.C SPRITES.C SOUND.C
\end{type}

\noindent For Pascal it is \texttt{MAIN=GAME.PAS} instead of
\texttt{SRC=} --- the units come in by \verb|uses|, from the same folder
--- and \texttt{OUT=} names the program when \texttt{NAME} in lower case
will not do. With a \texttt{PROJECT.K4P} beside the file in front, F9
builds the project, whichever of its files you are in --- a header or a
unit too --- and Ctrl-F9 runs the project's program; the files row starts
with its name in brackets. \verb|PROG GAME/PROJECT.K4P| (or
\emph{Open} on it) opens every source in a tab, and \emph{File}
$\rightarrow$ \emph{New project} makes the folder, the project file and a
first file that compiles as it stands. \verb|CC -p| and \verb|PAS -p|
build a project from the prompt the same way.

A C project compiles each source once and keeps what it made (on the
Linux side, never in your folder): the next F9 compiles only the files
that changed --- all of them if a \texttt{.H} beside them did --- and says
so, \texttt{game.prg: 4211 bytes (1 compiled, 2 kept)}.

\begin{aside}
\textbf{Where it is going.} This is PROG's fourth stage: projects came
third, the mouse and selection fourth, and REXX is already in. Next come
the other interpreters: EhBASIC, Microsoft BASIC, LOGO and Forth run on
the file in front of you, their errors in the same list.
\end{aside}

\section{Editing from inside a BASIC}\label{sec:editfrombasic}
There are two cases, and they look alike, so here is the rule.

\textbf{To edit the program you are writing}, type \verb|*EDIT| or
\verb|*VI| with nothing after it. BASIC saves the program to a temporary
file, runs the editor on it, and loads it back when you leave --- so what
you type in the editor is what you \verb|LIST| afterwards. Variables do
not survive the round trip, exactly as with \verb|LOAD|.
Section~\ref{sec:basicvi} has the details for EhBASIC; Microsoft BASIC
does the same with \verb|*VI| (Chapter~\ref{cha:msbasic}), BBC BASIC with
\verb|*VI| and \verb|*EDIT| (Chapter~\ref{cha:tube}), and LOGO with
\verb|EDIT "name| (Chapter~\ref{cha:logo}).

\textbf{To renumber it}, which none of those BASICs could do for itself,
use \verb|:renum| in VI or Ctrl-R in EDIT. Both know the language from the
file: a \texttt{.BAS} is EhBASIC's or Microsoft's, a \texttt{.BBC} is BBC
BASIC's, with its \verb|ELSE| targets and its lower-case variables (a
\texttt{goto} there is a name, not a jump). The lines' own numbers change,
and so does every target after \verb|GOTO|, \verb|GOSUB|, \verb|THEN|,
\verb|RESTORE| and \verb|ON|\ldots\verb|GOTO|; strings, \verb|REM| and
\verb|DATA| are left alone. A \verb|GOTO| to a line that does not exist is
kept as it is and counted in the message, and a program whose numbers are
out of order is refused, not guessed at. LOGO has no line numbers, and says
so.

\textbf{To edit any other file}, put \verb|SWAP| in front:

\begin{type}
*SWAP EDIT NOTES.TXT
\end{type}

\noindent A program loaded from the shell lands on top of whatever is in
memory, so a plain \verb|*EDIT NOTES.TXT| would run the editor over the
program you were writing. \verb|SWAP| (Chapter~\ref{cha:shell}) puts the
machine away first --- program, variables, screen and all --- runs the
editor on a clean one, and gives yours back when the editor exits.

So: no file name, no \verb|SWAP| needed, it is done for you; a file name,
\verb|SWAP| first.
